% from ega1-3-fr.tex
\oldpage[I]{111}
Il en résulte que pour toute extension $S'\to S$ du préschéma de base,
\phantomsection\label{I.3.3.12.base-change-monomorphism-fr}
\[
  f_{(S')}:X_{(S')}\to Y_{(S')}
\]
est un monomorphisme.
\end{env}

\begin{env}[3.3.13]
\label{I.3.3.13-fr}
Soient $S$, $S'$ deux schémas affines d'anneaux respectifs $A$, $A'$~; un
morphisme $S'\to S$ correspond donc à un homomorphisme d'anneaux $A\to A'$.
Si $X$ est un $S$-préschéma, on notera alors aussi $X_{(A')}$ ou
$X\otimes_A A'$ le $S'$-préschéma $X_{(S')}$~; lorsque $X$ est lui-même
affine d'anneau $B$, $X_{(A')}$ est affine d'anneau
$B_{(A')}=B\otimes_A A'$ obtenu par extension à $A'$ de l'anneau des scalaires
de la $A$-algèbre $B$.
\end{env}

\begin{env}[3.3.14]
\label{I.3.3.14-fr}
Avec les notations de (\hyperref[I.3.3.6-fr]{3.3.6}), pour tout
\emph{$S$-morphisme} $f:S'\to X$, $f'=(f,1_{S'})_S$ est un $S'$-morphisme
$S'\to X'=X_{(S')}$ tel que $p\circ f'=f$, $\pi'\circ f'=1_{S'}$, autrement
dit une \emph{$S'$-section de $X'$}~; et réciproquement si $f'$ est une telle
$S'$-section, $f=p\circ f'$ est un $S$-morphisme $S'\to X$. On définit donc
ainsi une \emph{correspondance biunivoque} canonique
\phantomsection\label{I.3.3.14.correspondence-fr}
\[
  \operatorname{Hom}_S(S',X)\overset{\sim}{\longrightarrow}
  \operatorname{Hom}_{S'}(S',X').
\]
On dit que $f'$ est le \emph{morphisme graphe} de $f$, et on le désigne par
$\Gamma_f$.
\end{env}

\begin{env}[3.3.15]
\label{I.3.3.15-fr}
Étant donné un préschéma $X$, qu'on peut toujours considérer comme un
$\mathbf{Z}$-préschéma, il résulte en particulier de
(\hyperref[I.3.3.14-fr]{3.3.14}) que les $X$-sections de
$X\otimes_{\mathbf{Z}}\mathbf{Z}[T]$ (où $T$ est une indéterminée)
correspondent biunivoquement aux \emph{morphismes}
$\mathbf{Z}[T]\to X$. Montrons que ces $X$-sections correspondent
biunivoquement aussi aux \emph{sections du faisceau structural
$\mathcal{O}_X$ au-dessus de $X$}. En effet, soit $(U_\alpha)$ un recouvrement
de $X$ par des ouverts affines~; soit
$u:X\to X\otimes_{\mathbf{Z}}\mathbf{Z}[T]$ un $X$-morphisme et soit
$u_\alpha$ sa restriction à $U_\alpha$~; si $A_\alpha$ est l'anneau du schéma
affine $U_\alpha$, $U_\alpha\otimes_{\mathbf{Z}}\mathbf{Z}[T]$ est un schéma
affine d'anneau $A_\alpha[T]$ (\hyperref[I.3.2.2-fr]{3.2.2}), et $u_\alpha$
correspond canoniquement à un $A_\alpha$-homomorphisme
$A_\alpha[T]\to A_\alpha$ (1.7.3). Or, un tel homomorphisme est complètement
déterminé par la donnée de l'image de $T$ dans $A_\alpha$, soit
$s_\alpha\in A_\alpha=\Gamma(U_\alpha,\mathcal{O}_X)$, et si on écrit que les
restrictions de $u_\alpha$ et de $u_\beta$ à un ouvert affine
$V\subset U_\alpha\cap U_\beta$ coïncident, on voit aussitôt que $s_\alpha$
et $s_\beta$ coïncident dans $V$~; donc la famille $(s_\alpha)$ est formée des
restrictions aux $U_\alpha$ d'une section $s$ de $\mathcal{O}_X$ au-dessus de
$X$~; et réciproquement, il est clair qu'une telle section définit une famille
$(u_\alpha)$ de morphismes qui sont les restrictions aux $U_\alpha$ d'un
$X$-morphisme $X\to X\otimes_{\mathbf{Z}}\mathbf{Z}[T]$. Ce résultat sera
généralisé dans II, 1.7.12.
\end{env}

\subsection{Points d'un préschéma à valeurs dans un préschéma~; points
géométriques}
\label{subsection:I.3.4-fr}

\begin{env}[3.4.1]
\label{I.3.4.1-fr}
Soit $X$ un préschéma~; pour tout préschéma $T$, on note encore $X(T)$
l'ensemble $\operatorname{Hom}(T,X)$ des morphismes $T\to X$, et les éléments de cet ensemble
sont aussi appelés \emph{points de $X$ à valeurs dans $T$}. Si on associe à
tout morphisme $f:T\to T'$ l'application $u'\mapsto u'\circ f$ de $X(T')$ dans
$X(T)$, on voit que, pour $X$ fixé, $X(T)$ est un \emph{foncteur contravariant
en $T$}, de la catégorie des préschémas dans celle des ensembles. En outre,
tout morphisme de préschémas $g:X\to Y$ définit un homomorphisme fonctoriel
$X(T)\to Y(T)$, faisant correspondre $g\circ v$ à $v\in X(T)$.
\end{env}

\begin{env}[3.4.2]
\label{I.3.4.2-fr}
Étant donnés trois ensembles $P$, $Q$, $R$ et deux applications
$\varphi:P\to R$, $\psi:Q\to R$, on appelle \emph{produit fibré de $P$ et $Q$
au-dessus de $R$} (relatif à $\varphi$ et $\psi$) la partie de
